\section{Conclusion}\label{sec:conclusion}

\subsection{The claim this paper actually supports}
This paper set out to answer one question honestly, in either direction:
does leakage-free graph topology add predictive value over a strong
transaction-only baseline? The answer, on the evidence gathered across four
datasets under one shared, leakage-audited protocol, is conditional rather
than universal, and we state it exactly as measured rather than rounded
toward a cleaner story. On BankSim, topology helps decisively: a
$+0.0974$ PR-AUC full-model structure lift ($11.7\%$ relative over a $0.833$
baseline), independently confirmed by a clean, amount-blind probe that still
finds $+0.1899$ PR-AUC once every amount feature is stripped from both sides
of the comparison. On the project's own synth\_erp corpus, topology helps by
the largest margin measured in this study ($+0.4031$, $246.8\%$ relative) --
but this number is explicitly barred from standing as thesis evidence by the
project's own policy, since the corpus was built to contain the structure
being tested for, and stands only as a demonstration of what topology
\emph{can} do under favorable conditions. On IBM AML, the comparison that
matters is uninformative rather than negative: an already-saturated $0.990$
baseline leaves almost no headroom, the nominal full-model lift
($+0.00792$) sits below this project's own noise floor, and the only honest
positive signal is a narrower, amount-blind probe ($+0.0905$) that isolates a
real but modest structural contribution once the dominant amount signal is
removed from both arms. On SAP W\"urzburg, topology does not help --
$-0.08288$ PR-AUC, a genuinely negative result on this project's real
general-ledger export, reported plainly rather than explained away, with the
caveat that only 25 test-split frauds make this specific number imprecise
even as its direction is corroborated by an independent, differently-built
experiment: a 46\% relative standard deviation on the 3-seed champion PR-AUC
for the same dataset (\S\ref{sec:limitations-stats}).
Four datasets, two clear directions and two qualified ones: this is the
paper's actual result, and it is deliberately not tidier than that.

We regard this qualified answer as more valuable than an inflated universal
one would have been, for a specific, falsifiable reason: this project's own
predecessor codebase claimed graph topology as ``the differentiator'' on
every dataset it touched, and that claim collapsed on inspection into $0.0$
feature importance and leakage-inflated metrics once anyone checked. A
four-for-four sweep of positive results, produced by the same team that
inherited that history, would be the least trustworthy possible outcome of
this study, not the most impressive one. The graph-based fraud detection
literature this paper positions itself within has, with some exceptions,
tended to report uniformly positive results for structural
features~\cite{motie2024gnn,yuan2025survey}; a controlled, same-architecture,
leakage-audited ablation that instead finds two positives, one null, and one
negative result -- and says so, including naming which of its own numbers
fall inside its own stated noise floor -- is offering a more falsifiable, and
therefore more useful, claim than a paper that reports success everywhere it
looked.

\subsection{A second, independent contribution: production engineering}
The paper's second contribution does not depend on the ablation's outcome at
all: a Procure-to-Pay ERP core with a forensic audit console was designed,
built, and \emph{measured}, not merely specified, against the hardening
requirements a real deployment would need. The online scoring path reuses
the research pipeline's own \texttt{WindowGraph} and \texttt{LapMemory}
classes directly rather than reimplementing them, verified by a parity test
that replays 20{,}000 real transactions through both the batch and online
computation paths with zero mismatches across all thirteen topology columns
(\S\ref{sec:system-arch}). Financial-integrity invariants -- double-entry
balance, fiscal-period close, and audit-log immutability -- are enforced as
Postgres triggers the application layer cannot bypass, not as conventions a
future migration could quietly break, and a five-layer duplicate-payment
defense was proven, not asserted, under genuine concurrency: 30 simultaneous
identical payment requests against one invoice produced exactly one
successful payment and 29 correctly rejected duplicates
(\S\ref{sec:system-integrity}). Performance and failure behavior are reported
as measured, including where the measurement is unflattering: the mixed
workload ramp clearly degrades before 500 concurrent users on a single
dev-mode process (p95 $8.18$s against a 1s target), and an initial
backup/restore drill correctly surfaced a Postgres MVCC-driven fingerprint
mismatch under concurrent writes rather than hiding it. None of this is
contingent on topology winning its ablation anywhere; it is a claim about
engineering discipline -- that a leakage-free feature computation can be
carried, unmodified, from a research pipeline into a live scoring path, and
that the resulting system's integrity guarantees, failure modes, and
performance ceiling can be demonstrated under real load rather than assumed.
We consider a genuinely deployed, tested, database-enforced system a
different and complementary kind of evidence from a PR-AUC number, and this
paper reports both rather than treating the console as an illustrative demo
of the research result.

\subsection{Future work}
Three directions follow directly from specific, named gaps this build
surfaced, rather than from general aspiration. First, and most concretely
scoped: persist the TGN's per-account GRU memory tensor as a checkpointed,
incrementally-updatable artifact, so that a live transaction can update and
reuse account memory in $O(1)$ rather than requiring a full history replay.
This is the single change that would let a true online champion -- topology
plus a learned temporal representation -- be served at all; today it is
blocked by an engineering gap (no persisted memory), not by an algorithmic
one, and closing it would also let the online/offline gap quantified in
Table~\ref{tab:online-gap} ($0.005$--$0.208$ PR-AUC by dataset) be measured
directly rather than only bounded. Second, extend multi-seed evaluation, and
ideally formal interval estimation, to the ablation itself rather than only
to the champion selection step: the seven-arm decomposition that produces
this paper's headline lift numbers remains single-seed, and work on
uncertainty-aware evaluation for graph-based fraud
models~\cite{shi2026uncertainty} points toward replacing this project's
informal $\sim$0.01 noise-floor heuristic with a proper bootstrap or paired
confidence interval, particularly for SAP W\"urzburg, where a 46\% relative
standard deviation in a related multi-seed experiment suggests the ablation's
point estimate alone should not be over-trusted. Third, resolve the batch
pipeline's z-score instability upstream rather than leaving it as a disclosed
but uncorrected defect: replacing the two-pass variance formula with
Welford's algorithm -- already done in the online engine -- would remove the
one place in this study where the offline research numbers and the online
serving numbers are known to diverge by construction, at the cost of
recomputing every headline PR-AUC this paper reports. We view that
recomputation, not the current disclosure-only stance, as the right eventual
resting state, and name it as such rather than treating disclosure as a
substitute for the fix.

The narrower, more durable claim this paper defends is not that graph
topology is or is not valuable for fraud detection in general, but that the
question can be asked honestly, with leakage structurally excluded by
construction rather than checked for after the fact, and that the answer one
gets is worth reporting exactly as measured -- including when it says no.
